% Point to the main file
\documentclass[../../Main.tex]{subfiles}

% ========== Start the document ==========
\begin{document}

\section{Section 4.2 - Binary}

Let's compute

$$67^2025 \pmod{77}$$

\begin{work}

    Let' see if FLT works

    \begin{center}
        $77 = 7 \cdot 11$ is \textbf{NOT} prime, so FLT does not apply
    \end{center}

    Let's see if EFLT works

    \begin{center}
        \begin{enumerate}
            \item $77 = 7 \cdot 11$ is the product of 2 distinct primes
            \item gcd$(67, 77) = 1$
        \end{enumerate}

        So EFLT applies
    \end{center}

    \begin{align*}
        6^{(7 - 1)(11 - 1)} &\equiv 1 \pmod{77} \\
        67^{60} &\equiv 1 \pmod{77} \\
    \end{align*}

    Breaking up 2025...
    $$2025 = 33 \cdot 60 + 45$$

    \begin{align*}
        67^{2025} \equiv 67^{45} \pmod{77} \\
    \end{align*}

    Break up 45 into powers of 2

    \begin{align*}
        67 \pmod{77} &= 67 \\
        67^2 \pmod{77} &= 23 \\
        67^4 \pmod{77} &= 23^2 \pmod{77} \\
        &= 67 \\
        67^8 \pmod{77} &= 23 \\
        67^{16} \pmod{77} &= 67 \\
        67^{32} \pmod{77} &= 23 \\
    \end{align*}

    \begin{align*}
        67^{2025} &\equiv 67^{45} \pmod{77} \\
        &\equiv 67^{32 + 8 + 4 + 1} \pmod{77} \\
        &\equiv 67^{32} \cdot 67^8 \cdot 67^4 \cdot 67^1 \pmod{77} \\
        &\equiv 23 \cdot 23 \cdot 67 \cdot 67 \cdot 67 \\
        &\equiv 23^2 \cdot 67^2 \pmod{77} \\
        \Aboxed{&\equiv 1} \pmod{77} \\
    \end{align*}

    How do we express 2025 as powers of 2?

    Let's list the powers of 2

    \begin{align*}
        2^0 &= 1 \\
        2^1 &= 2 \\
        2^2 &= 4 \\
        2^3 &= 8 \\
        2^4 &= 16 \\
        2^5 &= 32 \\
        &\vdots \\
        2^{10} &= 1024 \\
        2^{11} &= 2048 \\
    \end{align*}

    Now let's break up 2025 using a method similar to the Euclidean Algorithm

    \begin{align*}
        2025 &= 2^{10} + 1001 \\
        1001 &= 2^9 + 489 \\
        489 &= 2^8 + 233 \\
        233 &= 2^7 + 105 \\
        105 &= 2^6 + 41 \\
        41 &= 2^5 + 9 \\
        9 &= 2^3 + 2^0 \\
    \end{align*}

    $$2025 = 2^{10} + 2^9 + 2^8 + 2^7 + 2^6 + 2^5 + 2^3 + 2^0$$
    
    We can use this to write it in \textbf{binary representation}.

    \begin{align*}
        \underbrace{1}_{2^{10}}
        \underbrace{1}_{2^9}
        \underbrace{1}_{2^8}
        \underbrace{1}_{2^7}
        \underbrace{1}_{2^6}
        \underbrace{1}_{2^5}
        \underbrace{0}_{2^4}
        \underbrace{1}_{2^3}
        \underbrace{0}_{2^2}
        \underbrace{0}_{2^1}
        \underbrace{1_2}_{2^0}
    \end{align*}

    \begin{note}
        The 2 subscript denotes that this is in binary / base 2
    \end{note}

    Where 1 and 0 represent the coefficient attached to each power. Since it's 1 and 0, it simply represents the power being present and 0 being absent

    But notice that if the number $n$ is odd, then the last number of the binary representation is 1

    Else if $n$ is even then the last digit will be 0

    In other words

    \begin{center}
        $n \mod{2} = 1$ if odd \\
        $n \mod{2} = 0$ if even \\
    \end{center}

    \begin{align*}
        2025 &= 2 \cdot 1012 + \yellowbox{1} \\
        1012 &= 2 \cdot 506 + \yellowbox{0} \\
        506 &= 2 \cdot 253 + \yellowbox{0} \\
        253 &= 2 \cdot 126 + \yellowbox{1} \\
        126 &= 2 \cdot 63 + \yellowbox{0} \\
        63 &= 2 \cdot 31 + \yellowbox{1} \\
        31 &= 2 \cdot 15 + \yellowbox{1} \\
        15 &= 2 \cdot 7 + \yellowbox{1} \\
        7 &= 2 \cdot 3 + \yellowbox{1} \\
        3 &= 2 \cdot 1 + \yellowbox{1} \\
        1 &= 2 \cdot 0 + \yellowbox{1} \\
    \end{align*}

    Writing the sequence of remainders backwards we see

    \begin{align*}
        \underbrace{1}_{\text{last}}
        1 \;
        1 \;
        1 \;
        1 \;
        1 \;
        0 \;
        1 \;
        0 \;
        0 \;
        \underbrace{1_2}_{\text{first}}
        = 2025
    \end{align*}

\end{work}

But why does this work?

Given $n \in \setZ$

Want $n = a_k2^k + a_{k - 1}2^{k - 1} + \ldots + a_2 2^2 + a_1 2 + a_0$

Using the division algorithm,

$$n = 2q_0 - a_0$$

Then 

$$q_0 = \frac{n - a_0}{2}$$

Now repeat this with $q_0$

\begin{align*}
    q_0 = &2q_1 + a_1 \\
    &\vdots
\end{align*}

Let's do an example.

\begin{example}
    Write 1577 in binary

    \begin{align*}
        1577 &= 2 \cdot 788 + \yellowbox{1} \\
        788 &= 2 \cdot 394 + \yellowbox{0} \\
        394 &= 2 \cdot 197 + \yellowbox{0} \\
        197 &= 2 \cdot 98 + \yellowbox{1} \\
        98 &= 2 \cdot 49 + \yellowbox{0} \\
        49 &= 2 \cdot 24 + \yellowbox{1} \\
        24 &= 2 \cdot 12 + \yellowbox{0} \\
        12 &= 2 \cdot 6 + \yellowbox{0} \\
        6 &= 2 \cdot 3 + \yellowbox{0} \\
        3 &= 2 \cdot 1 + \yellowbox{1} \\
        1 &= 2 \cdot 0 + \yellowbox{1} \\
    \end{align*}

    Writing the bottom to top left to right...

    $$\boxed{11000101001}$$

    We could check this by just doing 

    $$1577 = 2^{10} + 2^9 + \ldots $$
    
\end{example}

\subsection{Base n Representation}

Y'know, 2 isn't special. We could actually use any number!

Let's try representing 1577 in base 8. Remember that each remainder we get is the coefficient to the corresponding power

\begin{align*}
    1577 &= 8 \cdot 197 + \yellowbox{1} \\
    197 &= 8 \cdot 24 + \yellowbox{5} \\
    24 &= 8 \cdot 3 + \yellowbox{0} \\
    3 &= 8 \cdot 0 + \yellowbox{3} \\
\end{align*}

Rewriting this, we get

\begin{align*}
    \underbrace{3}_{8^3}
    \underbrace{0}_{8^3}
    \underbrace{5}_{8^2}
    \underbrace{1}_{8^0}
\end{align*}

So 

$$1577 = 3 \cdot 8^3 + 5 \cdot 8^1 + 1 \cdot 8^0$$

Now let's do it with base 16

\begin{align*}
    1577 &= 16 \cdot 98 + \yellowbox{9} \\
    98 &= 16 \cdot 6 + \yellowbox{2} \\
    6 &= 16 \cdot 0 + \yellowbox{6} \\
\end{align*}

$$\boxed{1577 = 629_{16}}$$

$$1577 = 6 \cdot 16^2 + 2 \cdot 16^1 + 9 \cdot 16^0$$

For our last example, write the $AB0CD_{16}$ in base 10 (aka normal numbers)

Recall that

\begin{align*}
    A &= 10 \\
    B &= 11 \\
    &\vdots \\
\end{align*}

$AB0CD$ is just $10 \; 11 \; 0 \; 12 \; 13$

Since there's 5 "numbers" that means we start with $16^{5 - 1} = 16^4$

\begin{align*}
    AB0CD_{16} &= 10 \cdot 16^4 + 11 \cdot 16^3 + 0 \cdot 16^2 + 12 \cdot 16^1 + 13 \cdot 16^1 \\
    &= \boxed{700621} \tag{\text{just plug into calc}}
\end{align*}

% ========== End the document ==========
\end{document}